\Needspace{18\baselineskip}
\fancyhead[L]{\small GAUSS AN VON ZACH.}
\begin{center}
GAUSS AN VON ZACH.
\end{center}

\begin{center}
Monatliche Correspondenz. Band XIV. S. 187–192. August 1806.
\end{center}

\begin{flushright}
Braunschweig, 8. Juli 1806.
\end{flushright}

\par\indent --- --- --- Mithilfe der Mailänder Beobachtungen des Planeten \emph{Pallas} (M.\,C., Juliheft 1806, \href{https://www.digitale-sammlungen.de/en/view/bsb10538606?page=100}{S.~90}), die Sie mir freundlicherweise geschickt haben, habe ich neue Elemente (IX.) dieser Planetenbahn berechnet, die mit den Beobachtungen von \textsc{Oriani} möglichst genau übereinstimmen. Daraus habe ich die Opposition hergeleitet, was mit den zu stark abweichenden Elementen VIII nicht so gut möglich gewesen wäre. Das Ergebnis lautet:

\begin{center}
1805. Nov. 29.\ \ $11^{\mathrm{u}}\,18^{\mathrm{m}}\,17^{\mathrm{s}}$ mittlere Zeit in Seeberg.

Länge:\quad $67^\circ\,20'\,42''9$

Südliche geozentrische Breite:\quad $54\;\;30\;\;54.9$
\end{center}

\par\indent Für eine Ephemeride zur nächsten Erscheinung der \emph{Pallas} habe ich meine neuen Elemente nun auf die drei Oppositionen von 1803, 1804, 1805 und die Beobachtungen von 1802 gestützt. Dabei zeigte sich, dass sich diese verschiedenen Positionen nicht mehr ganz genau durch eine reine Ellipse darstellen ließen. Das ist zweifellos eine Folge der bisher vernachlässigten Störungen. Die Abweichung bleibt aber gering und dürfte kaum eine Minute überschreiten. Insbesondere ergibt sich aus den Beobachtungen von 1803 eine um $44''$ größere, aus denen von 1805 eine um ebenso viel kleinere Bahnneigung als der für die folgenden Elemente angesetzte Mittelwert. Der Knoten scheint nach den Beobachtungen von 1802 etwa $45''$ hinter dem angesetzten Mittel zu liegen, nach denen von 1804 nur ebenso weit davor. Hier sind die neuen Elemente (IX.), nach denen die folgende Ephemeride berechnet worden ist:
\Needspace{19\baselineskip}
\begin{center}\begingroup
\fontsize{9}{11}\selectfont
\setlength{\tabcolsep}{2pt}\renewcommand{\arraystretch}{1.12}
\begin{tabular}{lr}
\multicolumn{2}{c}{Epoche im Meridian von Seeberg} \\
1802. & $143^\circ\,24'\,1''1$ \\
1803. & $221\;32\;56.4$ \\
1804. & $299\;54\;42.5$ \\
1805. & $18\;3\;37.8$ \\
1806. & $96\;12\;33.2$ \\
1807. & $174\;21\;18.5$ \\
\end{tabular}\endgroup\end{center}
\begin{center}\begingroup
\fontsize{9}{11}\selectfont
\setlength{\tabcolsep}{2pt}\renewcommand{\arraystretch}{1.12}
\begin{tabular}{lr}
tägliche mittlere tropische Bewegung .\ .\ . & $770''7816$ \\
Länge des Aphels 1805 .\ .\ .\ . & $301^\circ\,11'\,31''1$ \\
Länge des aufsteigenden Knotens 1805 . & $172\;30\;5.8$ \\
Exzentrizität .\ .\ .\ .\ .\ .\ .\ .\ .\ .\ . & $0.2453840$ \\
Logarithmus der mittleren Entfernung .\ .\ . & $0.4421021$ \\
Neigung der Bahn .\ .\ .\ .\ .\ .\ .\ . & $34^\circ\,37'\,8''1$ \\
\end{tabular}\endgroup\end{center}

\par\indent Nach diesen Elementen würde die nächste Opposition am 4. Mai 1807 um 21 Uhr bei $223^\circ\,53'$ Länge und $42^\circ\,18'$ nördlicher geozentrischer Breite eintreten. Die Zeit wird zeigen, inwieweit diese Rechnung zutrifft. Meine Ephemeride habe ich gleich an Herrn Prof. \textsc{Harding} geschickt. Er wird danach eine Karte zeichnen, und ich hoffe, Sie werden uns bald in der M.\,C. damit beschenken.
\Needspace{15\baselineskip}
\SourceLeaf{283}%
\par\indent Meine nächste Arbeit wird nun \emph{Juno} betreffen. An meinen Elementen V wird wenig zu verbessern sein. Ich habe auch diese Rechnung bereits begonnen. Hier ist vorläufig mein Vergleich der Elemente V mit den letzten drei Beobachtungen, die Herr \textsc{Bessel} im Mai angestellt hat:

\begin{center}\begingroup
\fontsize{9}{11}\selectfont
\setlength{\tabcolsep}{2pt}\renewcommand{\arraystretch}{1.12}
\begin{tabular}{l|l|l|l|l}
& \multicolumn{2}{c|}{Berechnete} & \multicolumn{2}{c}{Unterschied in} \\
\shortstack{Mittlere Zeit\\in Lilienthal} & Rektaszension & \shortstack{Nördliche\\Deklination} & Rektaszension & Deklination \\
\hline
1806. Mai 17.\ $11^{\mathrm{u}}\,47^{\mathrm{m}}\,59^{\mathrm{s}}$ & $165^\circ\;4'\;11''$ & $9^\circ\;30'\;36''$ & $+0'\,51''$ & $+38''$ \\
19.\ $11\;\;1\;\;38$ & $165\;\;16\;\;34$ & $9\;\;29\;\;23$ & $+0\;\;33$ & $-4$ \\
23.\ $11\;\;31\;\;\;2$ & $165\;\;44\;\;57$ & $9\;\;25\;\;13$ & $+1\;\;11$ & $0$ \\
\cline{3-5}
& & Mittel & $+52''$ & $+11''$ \\
\end{tabular}\endgroup\end{center}

\par\indent Sobald ich die Elemente der \emph{Juno} verbessert habe, werde ich auch für sie gleich eine Ephemeride berechnen und sie Ihnen und Herrn Prof. \textsc{Harding} mitteilen. Zuletzt werde ich dieselbe Arbeit noch für \emph{Ceres} vornehmen, da ich hoffe, dass bis dahin mehr Beobachtungen aus der Zeit um die Opposition bekannt sein werden als bisher. [Der geozentrische Lauf der Pallas nach den Elementen IX ist vom 7. Dezember 1806 bis zum 21. September 1807 in Abständen von drei Tagen jeweils für 12 Uhr, Mitternacht in Seeberg, angegeben. Rektaszension und Deklination sind in Minuten angegeben, der Abstand von der Erde auf drei und die Lichtstärke auf vier Dezimalstellen.]
